\begin{table}[H]
\centering
\caption{Standardized Effect Sizes for Main Outcomes}
\label{tab:sde}
\begin{threeparttable}
\begin{tabular}{llcccccc}
\toprule
Outcome & Specification & $\hat{\beta}$ & SE & SD($Y$) & SDE & SE(SDE) & Classification \\
\midrule
\multicolumn{8}{l}{\textit{Panel A: Pooled}} \\
Log HPI & Full sample & 0.3108 & 0.0466 & 0.229 & 1.357 & 0.203 & Large positive \\
\midrule
\multicolumn{8}{l}{\textit{Panel B: Heterogeneous}} \\
Log HPI & Southern EU & 0.3683 & 0.0804 & 0.213 & 1.726 & 0.377 & Large positive \\
Log HPI & Northern/W. EU & 0.3307 & 0.0645 & 0.247 & 1.341 & 0.262 & Large positive \\
\bottomrule
\end{tabular}
\begin{tablenotes}[flushleft]
\small
\item \textit{Notes:} \textbf{Country:} Portugal (treated) versus 15 EU member state controls. \textbf{Research question:} Whether the abrupt termination of a preferential tax regime for foreign residents reduced residential house prices in the treated country relative to untreated EU peers. \textbf{Policy mechanism:} The Non-Habitual Resident regime granted new foreign tax residents a flat 20\% income tax rate and exemption on most foreign-sourced income for 10 years, attracting over 74,000 beneficiaries who concentrated in premium housing markets; its abolition removed the tax incentive for new arrivals and signaled regime instability to existing beneficiaries. \textbf{Outcome definition:} Log Eurostat House Price Index (prc\_hpi\_q), base 2015 = 100, measuring quarterly residential property transaction prices across all purchase types. \textbf{Treatment:} Binary; Portugal after the NHR termination announcement (2023Q3) versus all other periods and countries. \textbf{Data:} Eurostat quarterly HPI, 2010Q1--2025Q3, 16 EU countries, 945 country-quarter observations. \textbf{Method:} Two-way fixed effects difference-in-differences with country and quarter fixed effects; standard errors clustered at the country level. \textbf{Sample:} EU member states with continuous quarterly HPI coverage from 2010; excludes non-EU countries and states with missing data gaps. SDE $= \hat{\beta} / \text{SD}(Y)$ where SD($Y$) is the unconditional standard deviation of log HPI across the full sample. Classification refers to magnitude, not statistical significance: Large ($|$SDE$| > 0.15$), Moderate ($0.05$--$0.15$), Small ($0.005$--$0.05$), Null ($< 0.005$).
\end{tablenotes}
\end{threeparttable}
\end{table}

